\documentclass[12pt,letterpaper]{article}
\usepackage[margin=0.85in]{geometry}
\usepackage{amsmath,amssymb}
\usepackage[T1]{fontenc}
\setlength{\parindent}{0pt}
\setlength{\parskip}{0.7em}
\pagestyle{plain}
\newcommand{\question}[2]{\noindent\textbf{#1.}\hfill\textit{#2}\par\medskip}
\begin{document}
\begin{titlepage}
\thispagestyle{empty}
\vspace*{0.6in}
{\Large PHYS 200\par}
\vspace{0.25in}
{\LARGE\bfseries Practice Set 4\par}
\vspace{0.2in}
Cumulative review --- 8 questions\\
1 November 2026
\vspace{0.4in}

Six course questions and two generated questions.

Scope: relativity through light quanta, matter waves and two-slit interference, based on recorded progress through 1 November 2026.

Suggested time: 240 minutes, including double time. This is an optional practice estimate, not an official exam limit.

\textbf{Instructions}
\begin{itemize}
\item Answer all eight questions. Questions 1, 2 and 8 continue onto a second page. Use extra paper if needed.
\item Show your reasoning. Correct symbolic expressions or a sufficient set of equations determining the answer are acceptable; numerical arithmetic is optional.
\item Define frames and direction conventions, label diagrams, and include units. Give explanations where requested.
\item You may use your usual course formula sheet and calculator. Attempt Question 1 without the sheet first if you want an uncued terminology check.
\item No solutions are included. Course questions retain their givens and tasks, with selected subparts and minor typographical corrections. Question 4 retains its original part labels.
\end{itemize}
\vfill
Name: \rule{3.5in}{0.4pt}\par
Date: \rule{1.4in}{0.4pt}\hfill Start: \rule{0.8in}{0.4pt}\quad Finish: \rule{0.8in}{0.4pt}
\end{titlepage}
\setcounter{page}{1}


\question{1}{Generated}

Answer each part briefly, giving the physical meaning as well as any requested expression.


\textbf{(a)} State the two postulates of special relativity.


\vfill

\textbf{(b)} Define proper time and proper length. For each, state the measurement conditions that distinguish it from a quantity measured in another inertial frame.


\vfill

\textbf{(c)} Two events have separation \ensuremath{\Delta}x = 5 light-seconds and \ensuremath{\Delta}t = 3 s in one inertial frame. Classify their separation. Can an inertial frame make them simultaneous? Can one event causally influence the other?


\vfill

\newpage

\question{1 (continued)}{Generated}

\textbf{(d)} Identify the physical quantities represented by \(h/(mc)\), \(h/p\) and \(hc/\lambda\), stating what m, p and \ensuremath{\lambda} refer to in each case. How does the first quantity differ from the wavelength change in Compton scattering?


\vfill

\textbf{(e)} Explain the distinction between the sum of the constituent rest masses and the invariant mass of a system. Define positive binding energy relative to separated constituents and write the corresponding mass-energy relation.


\vfill

\textbf{(f)} What is the ultraviolet catastrophe? Name the law governing the wavelength of the peak of a blackbody spectrum. For a blackbody of fixed surface area whose temperature triples, determine the factors by which the peak wavelength, total emitted power and approximate photon emission rate change. For the last estimate, treat all emitted photons as having the peak wavelength.


\vfill

\newpage

\question{2}{Course}

The Avengers are battling the evil villain Thanos. Captain Marvel has attacked Thanos, and Thanos is now chasing her, both of them flying in the air. Doctor Strange is on the ground trying to help his team-mate. Captain Marvel flies past Doctor Strange at 0.6c, and 10 seconds later (in the Doctor's frame of reference) Thanos flies past him, travelling at 0.8c in the same direction as Captain Marvel.


\textbf{(a)} How fast is Thanos moving in Captain Marvel's frame of reference?


\vfill

\textbf{(b)} How far behind Captain Marvel was Thanos when Captain Marvel passed Doctor Strange, in the Doctor's frame? Give the answer in light-seconds.


\vfill

\textbf{(c)} If we now wanted to calculate how far Thanos was behind Captain Marvel at the moment that Captain Marvel passed Doctor Strange in the Captain's frame, then how many spacetime events do we need to consider?


\vfill

\textbf{(d)} Draw a spacetime diagram in Captain Marvel's frame, indicating the important spacetime events, the slopes of the relevant worldlines and the (unknown) distance ``X'' that Thanos was behind Captain Marvel when the Captain passed Doctor Strange in the Captain's frame.


\newpage

\question{2 (continued)}{Course}

\textbf{(e)} Hence (or otherwise), calculate the distance ``X'' that Thanos was behind Captain Marvel when the Captain passed Doctor Strange in the Captain's frame.


\vfill

\textbf{(f)} After Thanos passes Doctor Strange (but before he would have caught up with Captain Marvel), both the Doctor and the Captain fire energy bolts at Thanos (you should assume that these photonic beams and mystical energy bursts travel at the speed of light). In Thanos' frame of reference the two bursts are fired at the same time and they both arrive at the same (later) time. In Captain Marvel's frame, would the two blasts be fired at the same time and would they arrive at the same time?


\vfill

\textbf{(g)} Finally, considering Doctor Strange's frame, how long after Thanos passes Doctor Strange does Doctor Strange fire his mystic blast?


\newpage

\question{3}{Course}

A bullet, moving with speed 0.8c and with mass 5 g strikes a motionless toy car weighing 50 g. What is the speed of the toy car after the bullet wedges in it and how much mass does the toy car with the bullet wedged in it have?


\newpage

\question{4}{Course}

Batman, getting jealous of Superman's X-ray vision, decides that he wants UV vision for himself, and builds a photoelectric-effect-based gadget for measuring the wavelength of UV sources. During a calibration run, he discovers that the retarding potential required to stop electrons ejected by 3.07 \ensuremath{\times} 10\textsuperscript{15} Hz light is half that required to stop electrons ejected by 4.61 \ensuremath{\times} 10\textsuperscript{15} Hz light.


\textbf{(a)} What is the work function of the metal Batman uses?


\vfill

\textbf{(c)} Batman discovers that someone has planted a mystery UV source under his Batmobile. He points his new gadget at it and obtains a stopping voltage of 8 V for the photo-emitted electrons. What is the wavelength of the mystery source?


\newpage

\question{5}{Course}

\textbf{(a)} A photon of orange light, of wavelength 600 nm, scatters off a stationary electron, through an angle of 60\ensuremath{^\circ} relative to its original direction. By how much does the wavelength of the photon change, and what colour would the photon appear to be after it is Compton scattered?


\vfill

\textbf{(b)} What would be the biggest possible shift in the wavelength, for any scattering angle? Based on this, can you say whether Compton scattering is important for visible light?


\vfill

\textbf{(c)} Would Compton scattering for visible photons be larger or smaller if the photons scattered off protons rather than electrons?


\newpage

\question{6}{Course}

In this question, we will consider a Sci-Fi concept, a ship driven by an Earth-based laser. This works as follows: a really powerful and well collimated laser beam is sent from Earth in the direction the ship is supposed to go in. The spaceship is equipped with a large mirror and the laser beam is reflected back to the Earth. Since light carries momentum, this propels the ship forward.


Let's consider a time after the initial boost-off, when the laser is off and the ship is moving with a large velocity v away from the Earth. To give the ship's velocity a final adjustment, a single short laser burst is emitted towards the ship, with total energy E in the Earth's frame of reference. How much momentum, in the Earth's frame of reference, does this burst give to the ship once it's been reflected off its mirror?


You may assume that the spaceship is very massive and therefore, in the ship's frame of reference the momentum of the reflected pulse has the same magnitude as the momentum of the incident pulse. Equivalently, in the ship's frame of reference, the reflected light has the same wavelength as the incident light.


\newpage

\question{7}{Course}

In a certain classroom, the students notice that there are two holes in the back wall, separated by 1 m. They decide (for reasons known only to themselves) to see if it is possible to demonstrate the double-slit experiment by throwing their cell-phones at the two holes and looking for the interference pattern on a distant wall, 1 km away. If we assume that it's only possible to discern an interference pattern in the cell-phone positions if the fringes are at least 0.1 m apart on the wall (i.e. \ensuremath{\Delta}Y \ensuremath{\simeq} 0.1 m), then what is the required wavelength of the phones? If the average phone weighs 100 g, then what typical speed must they have in order for us to see this big an interference pattern? Finally, estimate how long it takes the phones to cover the 1 km distance. Is this a realistic experiment to perform?


\newpage

\question{8}{Generated}

A beam of free electrons, each with kinetic energy 4.5 eV in the laboratory, is directed at two narrow slits in a stationary barrier. The apparatus is in vacuum. The electrons are emitted one at a time, and the two slit contributions have equal amplitude and are in phase at the slits.


\textbf{(a)} Determine the electron wavelength in nm. State whether a nonrelativistic treatment is justified. You may use m\textsubscript{e} = 9.11 \ensuremath{\times} 10\textsuperscript{-31} kg, h = 6.626 \ensuremath{\times} 10\textsuperscript{-34} J s and 1 eV = 1.602 \ensuremath{\times} 10\textsuperscript{-19} J.


\vfill

\textbf{(b)} At a point P on the screen, the path from the lower slit is longer than the path from the upper slit by 3.5 electron wavelengths. Determine the phase difference modulo \(2\pi\) and whether P is a maximum or minimum. Sketch and label the two paths.


\vfill

\newpage

\question{8 (continued)}{Generated}

\textbf{(c)} Describe individual electron detections and the pattern built up after many electrons. Then describe what changes if a detector reliably records which slit each electron passes through. Explain the difference between a coherent superposition of the two paths and an incoherent mixture of the two alternatives.


\vfill

\textbf{(d)} A monochromatic light beam is now sent through the same stationary slits in vacuum. Does passing through the slits change its frequency or wavelength? Explain.


\vfill

\textbf{(e)} Keeping the same illuminated area and average intensity, the light frequency is doubled. What happens to the photon energy, photon momentum, photon arrival rate and electric-field amplitude? Justify each change.


\vfill

\end{document}